\documentclass[11pt]{scrartcl}

\usepackage[top=1.5cm]{geometry}
\usepackage{url}
\usepackage{float}
\usepackage{listings}
\usepackage{xcolor}
\usepackage{graphicx}
\usepackage{booktabs}

\setlength{\parindent}{0em}
\setlength{\parskip}{0.5em}

\newcommand{\youranswerhere}{[Your answer goes here \ldots]}
\renewcommand{\thesubsection}{\arabic{subsection}}

\lstdefinestyle{dmrsql}{
  language=SQL,
  basicstyle=\small\ttfamily,
  keywordstyle=\color{magenta!75!black},
  stringstyle=\color{green!50!black},
  showspaces=false,
  showstringspaces=false,
  commentstyle=\color{gray}}

\lstdefinestyle{dmrJava}{
  language=JAVA,
  basicstyle=\small\ttfamily,
  keywordstyle=\color{magenta!75!black},
  stringstyle=\color{green!50!black},
  showspaces=false,
  showstringspaces=false,
  commentstyle=\color{gray}}

\title{
  \textbf{\large Projektaufgabe 3 } \\
  Phase 3 – Implementierung des XPath Accelerators (3.5 P) \\
  {\large Datenmanagement jenseits von Relationen}
}

\author{
  Gruppe A08 \\
  \large Acikyol Aleyna, 12032843 \\
  \large Özcan Ahmet, 12311287 \\
  \large Sadykova Nargiz, 12129337
}

\begin{document}

\maketitle\thispagestyle{empty}

Dieses Reporting Template dient der Vorbereitung der Abgabe von Phase 3.

\subsection*{Verkleinern der Fensteranfrage (0.5 P)}

Zeigen Sie, dass die Berechnung der Achsen für die Beispiele aus Phase 1 hier die gleichen Ergebnisse erzeugen:

\begin{table}[h]
	\centering
		\begin{center}
			\begin{tabular}{ l | c c }
				\toprule
				Achse & Ergebnisknoten ID's & Ergebnisgröße\\
				\midrule
				ancestor & 1 bis 4 & 4 \\
				descendants & 4 bis 31 & 28 \\
				following SchmittKAMM23& 19 & 1 \\
				preceeding SchmittKAMM23 & - & 0 \\
				following SchalerHS23& - & 0 \\
				preceeding SchalerHS23& 4 & 1 \\
				\bottomrule
			\end{tabular}
			\end{center}
	\caption{Ergebnisse der XPath-Berechnung auf Edge Modell aus Phase 1}
	\label{tab:ErgebnisseDerXPathBerechnug}
\end{table}

\begin{table}[h]
	\centering
		\begin{center}
			\begin{tabular}{ l | c c }
				\toprule
				Achse & Ergebnisknoten ID's & Ergebnisgröße\\
				\midrule
				ancestor & 1 bis 4 & 4 \\
				descendants & 4 bis 31 & 28 \\
				following SchmittKAMM23& 19 & 1 \\
				preceeding SchmittKAMM23 & - & 0 \\
				following SchalerHS23& - & 0 \\
				preceeding SchalerHS23& 4 & 1 \\
				\bottomrule
			\end{tabular}
			\end{center}
	\caption{Ergebnisse der XPath-Berechnung des XPath-Accelerators aus Phase 3 mit Verkleinerung des Fensters}
	\label{tab:ErgebnisseDerXPathBerechnug}
\end{table}


\newpage
\subsection*{Schemaerstellung (1 Punkt).}
Geben Sie die Create Table statements aller von Ihnen angelegten Tabellen an:
\begin{lstlisting}[style=dmrsql]
CREATE TABLE accel (
  pre    INTEGER,
  post   INTEGER,
  parent INTEGER,
  kind   TEXT,
  name   TEXT
);

CREATE TABLE content (
  pre  INTEGER,
  text TEXT
);

CREATE TABLE attribute (
  pre  INTEGER,
  text TEXT
);

CREATE TABLE accel1d (
  pre    INTEGER,
  post   INTEGER,
  parent INTEGER,
  kind   TEXT,
  name   TEXT
);

\end{lstlisting}
\newpage
\subsection*{Zugriff mit nur einer Achse (1 Punkt).}
Zeigen Sie, dass die Berechnung der Achsen für die Beispiele aus Phase 1 hier die gleichen Ergebnisse erzeugen:

\begin{table}[h]
	\centering
		\begin{center}
			\begin{tabular}{ l | c c }
				\toprule
				Achse & Ergebnisknoten ID's & Ergebnisgröße\\
				\midrule
				descendants & 4 bis 57 & 28 \\
				\bottomrule
			\end{tabular}
			\end{center}
	\caption{Ergebnisse der XPath-Berechnung des XPath-Accelerators aus Phase 3 mit nur einer Achse}
	\label{tab:ErgebnisseDerXPathBerechnug}
\end{table}

Da sich die Knoten IDs unterscheiden, zeigen Sie für einige Knoten die jeweilige Identität.

\begin{itemize}
    \item  (11, 12, 4, 'element', 'author')
   \item(13, 14, 4, 'element', 'author')
   \item(15, 16, 4, 'element', 'title')
   \item(17, 18, 4, 'element', 'pages')
   \item(19, 20, 4, 'element', 'year')
   \item(21, 22, 4, 'element', 'volume')
\end{itemize}

\subsection*{Zeitmamagement}

Benötigte Zeit pro Person (nur Phase 3): \textbf{Nargiz Sadykova: 3 Std., Aleyna Acikyol: 5Std., }


\end{document}
